% assets/w07-rdot-rthetadot.svg — rdot 과 r*thetadot 이 그림의 어느 선인가
%
%   bash _tools/tikz2svg.sh assets/src/w07-rdot-rthetadot.tex assets/w07-rdot-rthetadot.svg
%
% 요점 둘
%   (1) 짧은 시간 dt 동안 배가 움직인 자취를 두 조각으로 쪼개면 그 둘이 바로
%       반경 조각(dr)과 호 조각(r*dtheta)이다. 두 조각은 서로 직각이다.
%   (2) 같은 둘을 속도에서 보면 — 속도 화살표를 두 축에 내린 그림자(정사영)다.
%
% 왼쪽 그림의 dtheta 는 보이라고 크게 그린 것이다. 실제로는 아주 작다.
% 호 길이 = 반경 x 각 (rad) 이므로 r*dtheta 이고, dt 로 나누면 r*thetadot 이다.

\documentclass[border=6pt]{standalone}
% capstone-style.tex — 이 볼트의 TikZ 그림이 공유하는 스타일
%
% 쓰는 법:  % capstone-style.tex — 이 볼트의 TikZ 그림이 공유하는 스타일
%
% 쓰는 법:  % capstone-style.tex — 이 볼트의 TikZ 그림이 공유하는 스타일
%
% 쓰는 법:  \input{capstone-style}   (assets/src/ 안에서)
% 빌드:     bash _tools/tikz2svg.sh assets/src/w02-node-topic.tex assets/w02-node-topic.svg
%
% 왜 TikZ 인가
%   손으로 SVG 좌표를 쓰면 화살촉·선 굵기·글자 크기가 그림마다 미묘하게 달라진다.
%   그것이 "자료가 어설퍼 보인다"의 정체다. 스타일을 한 번만 정하고 상대좌표로
%   배치하면 좌표를 고쳐도 그림이 무너지지 않는다.
%   수식은 본문과 같은 TeX 글꼴로 나온다.
%
% 한글
%   ko.TeX(DVI 경로)를 쓴다. XeLaTeX 은 TikZ 그래픽을 PDF special 로 내보내는데
%   dvisvgm 이 그것을 받으려면 Ghostscript < 10.01 이 필요하다. 이 PC 에는
%   10.07 이 깔려 있어 그 경로가 막힌다. latex -> DVI -> dvisvgm 이 유일하게
%   동작하는 조합이다.

\usepackage[hangul]{kotex}
\usepackage{tikz}
\usepackage{amsmath}
\usepackage{xcolor}
\usetikzlibrary{arrows.meta, positioning, calc, fit, backgrounds, decorations.pathmorphing}

% ---- 색 --------------------------------------------------------------------
% 강조는 하나만. 나머지는 검정. 흑백 인쇄에서도 읽혀야 한다.
\definecolor{ink}{HTML}{1A1A1A}
\definecolor{soft}{HTML}{666666}
\definecolor{accent}{HTML}{7C3AED}
\definecolor{warn}{HTML}{B91C1C}

% 계층 색 — Simulink 모델의 블록 색과 같은 규칙을 쓴다
\definecolor{guid}{HTML}{1D4ED8}    % 유도
\definecolor{ctrl}{HTML}{EA580C}    % 제어
\definecolor{thr}{HTML}{D97706}     % 추진기
\definecolor{plant}{HTML}{16A34A}   % 운동모델
\definecolor{nav}{HTML}{7C3AED}     % 항법 · 신호처리
\definecolor{pathc}{HTML}{0D9488}   % 계획 경로 · 항적
\definecolor{errc}{HTML}{C2410C}    % 오차
\definecolor{flow}{HTML}{0369A1}    % 조류 · 바람

% ---- 선과 화살촉 -----------------------------------------------------------
\tikzset{
  sig/.style   = {ink, line width=0.5pt, -{Stealth[length=4pt, width=3pt]}},
  wire/.style  = {ink, line width=0.5pt},
  fb/.style    = {ink, line width=0.5pt, densely dashed,
                  -{Stealth[length=4pt, width=3pt]}},
  asig/.style  = {accent, line width=0.5pt, densely dashed,
                  -{Stealth[length=4pt, width=3pt]}},
  % 블록 — 흰 바탕, 직각, 검은 테두리. 색은 테두리에만 준다
  blk/.style   = {draw=ink, line width=0.55pt, fill=white,
                  minimum width=18mm, minimum height=8.5mm, inner sep=2pt, font=\small},
  gblk/.style  = {blk, draw=guid},
  cblk/.style  = {blk, draw=ctrl},
  tblk/.style  = {blk, draw=thr},
  pblk/.style  = {blk, draw=plant},
  nblk/.style  = {blk, draw=nav},
  % 합산점 · 분기점
  sum/.style   = {draw=ink, line width=0.55pt, fill=white, circle,
                  inner sep=0pt, minimum size=4.6mm, font=\scriptsize},
  dot/.style   = {ink, fill=ink, circle, inner sep=0pt, minimum size=1.5mm},
  % 글자
  sname/.style = {font=\small, inner sep=1.5pt},
  note/.style  = {font=\scriptsize, soft, inner sep=1.5pt},
  wnote/.style = {font=\scriptsize, warn, inner sep=1.5pt},
  anote/.style = {font=\scriptsize, accent, inner sep=1.5pt},
  axis/.style  = {ink, line width=0.5pt},
  tick/.style  = {font=\scriptsize, soft},
  % 묶음 상자 — 서브시스템 표시
  grp/.style   = {draw=soft, line width=0.4pt, densely dotted, rounded corners=1mm,
                  inner sep=3mm},
}
   (assets/src/ 안에서)
% 빌드:     bash _tools/tikz2svg.sh assets/src/w02-node-topic.tex assets/w02-node-topic.svg
%
% 왜 TikZ 인가
%   손으로 SVG 좌표를 쓰면 화살촉·선 굵기·글자 크기가 그림마다 미묘하게 달라진다.
%   그것이 "자료가 어설퍼 보인다"의 정체다. 스타일을 한 번만 정하고 상대좌표로
%   배치하면 좌표를 고쳐도 그림이 무너지지 않는다.
%   수식은 본문과 같은 TeX 글꼴로 나온다.
%
% 한글
%   ko.TeX(DVI 경로)를 쓴다. XeLaTeX 은 TikZ 그래픽을 PDF special 로 내보내는데
%   dvisvgm 이 그것을 받으려면 Ghostscript < 10.01 이 필요하다. 이 PC 에는
%   10.07 이 깔려 있어 그 경로가 막힌다. latex -> DVI -> dvisvgm 이 유일하게
%   동작하는 조합이다.

\usepackage[hangul]{kotex}
\usepackage{tikz}
\usepackage{amsmath}
\usepackage{xcolor}
\usetikzlibrary{arrows.meta, positioning, calc, fit, backgrounds, decorations.pathmorphing}

% ---- 색 --------------------------------------------------------------------
% 강조는 하나만. 나머지는 검정. 흑백 인쇄에서도 읽혀야 한다.
\definecolor{ink}{HTML}{1A1A1A}
\definecolor{soft}{HTML}{666666}
\definecolor{accent}{HTML}{7C3AED}
\definecolor{warn}{HTML}{B91C1C}

% 계층 색 — Simulink 모델의 블록 색과 같은 규칙을 쓴다
\definecolor{guid}{HTML}{1D4ED8}    % 유도
\definecolor{ctrl}{HTML}{EA580C}    % 제어
\definecolor{thr}{HTML}{D97706}     % 추진기
\definecolor{plant}{HTML}{16A34A}   % 운동모델
\definecolor{nav}{HTML}{7C3AED}     % 항법 · 신호처리
\definecolor{pathc}{HTML}{0D9488}   % 계획 경로 · 항적
\definecolor{errc}{HTML}{C2410C}    % 오차
\definecolor{flow}{HTML}{0369A1}    % 조류 · 바람

% ---- 선과 화살촉 -----------------------------------------------------------
\tikzset{
  sig/.style   = {ink, line width=0.5pt, -{Stealth[length=4pt, width=3pt]}},
  wire/.style  = {ink, line width=0.5pt},
  fb/.style    = {ink, line width=0.5pt, densely dashed,
                  -{Stealth[length=4pt, width=3pt]}},
  asig/.style  = {accent, line width=0.5pt, densely dashed,
                  -{Stealth[length=4pt, width=3pt]}},
  % 블록 — 흰 바탕, 직각, 검은 테두리. 색은 테두리에만 준다
  blk/.style   = {draw=ink, line width=0.55pt, fill=white,
                  minimum width=18mm, minimum height=8.5mm, inner sep=2pt, font=\small},
  gblk/.style  = {blk, draw=guid},
  cblk/.style  = {blk, draw=ctrl},
  tblk/.style  = {blk, draw=thr},
  pblk/.style  = {blk, draw=plant},
  nblk/.style  = {blk, draw=nav},
  % 합산점 · 분기점
  sum/.style   = {draw=ink, line width=0.55pt, fill=white, circle,
                  inner sep=0pt, minimum size=4.6mm, font=\scriptsize},
  dot/.style   = {ink, fill=ink, circle, inner sep=0pt, minimum size=1.5mm},
  % 글자
  sname/.style = {font=\small, inner sep=1.5pt},
  note/.style  = {font=\scriptsize, soft, inner sep=1.5pt},
  wnote/.style = {font=\scriptsize, warn, inner sep=1.5pt},
  anote/.style = {font=\scriptsize, accent, inner sep=1.5pt},
  axis/.style  = {ink, line width=0.5pt},
  tick/.style  = {font=\scriptsize, soft},
  % 묶음 상자 — 서브시스템 표시
  grp/.style   = {draw=soft, line width=0.4pt, densely dotted, rounded corners=1mm,
                  inner sep=3mm},
}
   (assets/src/ 안에서)
% 빌드:     bash _tools/tikz2svg.sh assets/src/w02-node-topic.tex assets/w02-node-topic.svg
%
% 왜 TikZ 인가
%   손으로 SVG 좌표를 쓰면 화살촉·선 굵기·글자 크기가 그림마다 미묘하게 달라진다.
%   그것이 "자료가 어설퍼 보인다"의 정체다. 스타일을 한 번만 정하고 상대좌표로
%   배치하면 좌표를 고쳐도 그림이 무너지지 않는다.
%   수식은 본문과 같은 TeX 글꼴로 나온다.
%
% 한글
%   ko.TeX(DVI 경로)를 쓴다. XeLaTeX 은 TikZ 그래픽을 PDF special 로 내보내는데
%   dvisvgm 이 그것을 받으려면 Ghostscript < 10.01 이 필요하다. 이 PC 에는
%   10.07 이 깔려 있어 그 경로가 막힌다. latex -> DVI -> dvisvgm 이 유일하게
%   동작하는 조합이다.

\usepackage[hangul]{kotex}
\usepackage{tikz}
\usepackage{amsmath}
\usepackage{xcolor}
\usetikzlibrary{arrows.meta, positioning, calc, fit, backgrounds, decorations.pathmorphing}

% ---- 색 --------------------------------------------------------------------
% 강조는 하나만. 나머지는 검정. 흑백 인쇄에서도 읽혀야 한다.
\definecolor{ink}{HTML}{1A1A1A}
\definecolor{soft}{HTML}{666666}
\definecolor{accent}{HTML}{7C3AED}
\definecolor{warn}{HTML}{B91C1C}

% 계층 색 — Simulink 모델의 블록 색과 같은 규칙을 쓴다
\definecolor{guid}{HTML}{1D4ED8}    % 유도
\definecolor{ctrl}{HTML}{EA580C}    % 제어
\definecolor{thr}{HTML}{D97706}     % 추진기
\definecolor{plant}{HTML}{16A34A}   % 운동모델
\definecolor{nav}{HTML}{7C3AED}     % 항법 · 신호처리
\definecolor{pathc}{HTML}{0D9488}   % 계획 경로 · 항적
\definecolor{errc}{HTML}{C2410C}    % 오차
\definecolor{flow}{HTML}{0369A1}    % 조류 · 바람

% ---- 선과 화살촉 -----------------------------------------------------------
\tikzset{
  sig/.style   = {ink, line width=0.5pt, -{Stealth[length=4pt, width=3pt]}},
  wire/.style  = {ink, line width=0.5pt},
  fb/.style    = {ink, line width=0.5pt, densely dashed,
                  -{Stealth[length=4pt, width=3pt]}},
  asig/.style  = {accent, line width=0.5pt, densely dashed,
                  -{Stealth[length=4pt, width=3pt]}},
  % 블록 — 흰 바탕, 직각, 검은 테두리. 색은 테두리에만 준다
  blk/.style   = {draw=ink, line width=0.55pt, fill=white,
                  minimum width=18mm, minimum height=8.5mm, inner sep=2pt, font=\small},
  gblk/.style  = {blk, draw=guid},
  cblk/.style  = {blk, draw=ctrl},
  tblk/.style  = {blk, draw=thr},
  pblk/.style  = {blk, draw=plant},
  nblk/.style  = {blk, draw=nav},
  % 합산점 · 분기점
  sum/.style   = {draw=ink, line width=0.55pt, fill=white, circle,
                  inner sep=0pt, minimum size=4.6mm, font=\scriptsize},
  dot/.style   = {ink, fill=ink, circle, inner sep=0pt, minimum size=1.5mm},
  % 글자
  sname/.style = {font=\small, inner sep=1.5pt},
  note/.style  = {font=\scriptsize, soft, inner sep=1.5pt},
  wnote/.style = {font=\scriptsize, warn, inner sep=1.5pt},
  anote/.style = {font=\scriptsize, accent, inner sep=1.5pt},
  axis/.style  = {ink, line width=0.5pt},
  tick/.style  = {font=\scriptsize, soft},
  % 묶음 상자 — 서브시스템 표시
  grp/.style   = {draw=soft, line width=0.4pt, densely dotted, rounded corners=1mm,
                  inner sep=3mm},
}


\begin{document}
\begin{tikzpicture}[x=1mm, y=1mm]

% ===========================================================================
% 왼쪽 — 짧은 시간 동안 움직인 자취를 두 조각으로 쪼갠다
% ===========================================================================
\begin{scope}[shift={(-48,0)}]

  \def\Rone{42}      % 지금 거리 r
  \def\Rtwo{34}      % 조금 뒤 거리 r + dr  (dr < 0, 안으로 들어왔다)
  \def\Aone{72}      % 지금 위치의 TikZ 각
  \def\Atwo{44}      % 조금 뒤 위치의 TikZ 각 (시계로 돌았다)

  \coordinate (O)  at (0,0);
  \coordinate (P1) at (\Aone:\Rone mm);
  \coordinate (Q)  at (\Atwo:\Rone mm);
  \coordinate (P2) at (\Atwo:\Rtwo mm);

  \node[sname, anchor=south] at (6,50) {$\Delta t$ 동안 움직인 자취를 두 조각으로};

  % 중심과 두 반경선
  \fill[ink] (O) circle (1.2);
  \node[note, anchor=north east, inner sep=1pt] at (-1,-1) {중심};
  \draw[ink, line width=0.6pt] (O) -- (P1);
  \draw[ink, line width=0.6pt] (O) -- (Q);
  \node[note, fill=white, inner sep=0.8pt] at (\Aone:24mm) {$r$};
  \node[note, fill=white, inner sep=0.8pt] at (\Atwo:21mm) {$r+\Delta r$};

  % 중심의 dtheta
  \draw[errc, line width=0.6pt, -{Stealth[length=3pt,width=2.4pt]}]
        (\Aone:12mm) arc (\Aone:\Atwo:12mm);
  \node[note, errc, anchor=west, inner sep=1pt] at (\Atwo+14:14.5mm) {$\Delta\theta$};

  % 호 조각 — 반경을 그대로 둔 채 각만 바뀐 몫
  \draw[thr, line width=1.2pt, -{Stealth[length=4.5pt,width=3.5pt]}]
        (\Aone:\Rone mm) arc (\Aone:\Atwo:\Rone mm);
  \node[note, thr, anchor=south west, text width=30mm] at (\Aone-13:46mm)
       {호 조각 — 각만 바뀐 몫\\ 길이 $= r\,\Delta\theta$};

  % 반경 조각 — 각을 그대로 둔 채 거리만 바뀐 몫
  \draw[guid, line width=1.2pt, -{Stealth[length=4.5pt,width=3.5pt]}] (Q) -- (P2);
  \node[note, guid, anchor=west, text width=26mm] at ($(Q)+(3.0,1.0)$)
       {반경 조각 — 거리만 바뀐 몫\\ 길이 $= \Delta r$ (지금은 음수)};

  % 두 조각은 직각이다
  \draw[soft, line width=0.4pt]
        ($(Q)+(\Atwo+90:2.4mm)$) -- ($(Q)+(\Atwo+90:2.4mm)+(\Atwo+180:2.4mm)$)
        -- ($(Q)+(\Atwo+180:2.4mm)$);

  % 실제로 간 길 — 두 조각의 합
  \draw[plant, line width=0.7pt, densely dashed, -{Stealth[length=4pt,width=3pt]}]
        (P1) -- (P2);

  \fill[errc] (P1) circle (1.5);
  \node[sname, errc, anchor=south east, inner sep=1.5pt] at ($(P1)+(-0.6,0.8)$) {지금};
  \fill[plant] (P2) circle (1.5);
  \node[sname, plant, anchor=north west, inner sep=1.5pt] at ($(P2)+(0.8,-0.6)$)
       {$\Delta t$ 뒤};

  \node[note, anchor=north west, text width=58mm] at (-26,-14)
       {두 조각을 $\Delta t$ 로 나누면 그대로 속도 성분이 된다\\[3pt]
        \textcolor{guid}{$\dfrac{\Delta r}{\Delta t} \rightarrow \dot{r}$}
        \qquad
        \textcolor{thr}{$\dfrac{r\,\Delta\theta}{\Delta t} \rightarrow r\dot{\theta}$}};

  \node[note, anchor=north west, text width=52mm] at (-26,-33)
       {$\Delta\theta$ 는 보이라고 크게 그린 것이다. 실제로는 아주 작고,
        그때 호 조각은 직선과 구별되지 않는다};

\end{scope}

% ===========================================================================
% 오른쪽 — 같은 둘을 속도에서 보면 그림자다
% ===========================================================================
\begin{scope}[shift={(46,0)}]

  \def\Rp{40}        % 중심에서 배까지
  \def\Ap{58}        % 배 위치의 TikZ 각
  \def\Aer{58}       % e_r 방향 = 배 위치 각
  \def\Aet{-32}      % e_theta 방향 = 그보다 90도 시계

  \coordinate (O) at (0,0);
  \coordinate (P) at (\Ap:\Rp mm);

  \node[sname, anchor=south] at (2,50) {속도 화살표가 두 축에 드리우는 \textbf{그림자}};

  \fill[ink] (O) circle (1.2);
  \node[note, anchor=north east, inner sep=1pt] at (-1,-1) {중심};
  \draw[ink, line width=0.6pt, -{Stealth[length=4pt,width=3pt]}] (O) -- (P);
  \node[note, fill=white, inner sep=0.8pt] at (\Ap:22mm) {$r$};

  % 배 자리의 두 축 — e_r 은 반경선의 연장, e_theta 는 거기에 수직
  \draw[soft, line width=0.5pt, -{Stealth[length=3.5pt,width=2.8pt]}] (P) -- ++(\Aer:22mm);
  \node[note, anchor=south west, text width=24mm, inner sep=1pt] at ($(P)+(\Aer:22.5mm)$)
       {$\hat{e}_r$ — 중심에서 멀어지는 쪽};
  \draw[soft, line width=0.5pt, -{Stealth[length=3.5pt,width=2.8pt]}] (P) -- ++(\Aet:30mm);
  \node[note, anchor=north west, inner sep=1pt] at ($(P)+(\Aet:30.5mm)$)
       {$\hat{e}_\theta$ — 도는 쪽};
  % 두 축은 직각
  \draw[soft, line width=0.4pt]
        ($(P)+(\Aer:3mm)$) -- ($(P)+(\Aer:3mm)+(\Aet:3mm)$) -- ($(P)+(\Aet:3mm)$);

  % 속도 — 안으로 들어오며 도는 중
  % 각 = \Aer - alpha = 58 - 125 = -67. 두 그림자의 합과 정확히 맞춘 값이다
  \draw[plant, line width=1.4pt, -{Stealth[length=5pt,width=4pt]}] (P) -- ++(-67:27mm);
  \node[sname, plant, anchor=west, fill=white, inner sep=1pt] at ($(P)+(-67:28mm)$)
       {$v_d$};
  \coordinate (V) at ($(P)+(-67:27mm)$);

  % 그림자 둘 — 수선을 내린다
  \draw[soft, line width=0.35pt, densely dotted] (V) -- ($(P)+(\Aer+180:15.5mm)$);
  \draw[soft, line width=0.35pt, densely dotted] (V) -- ($(P)+(\Aet:22.1mm)$);

  \draw[guid, line width=1.2pt, -{Stealth[length=4.5pt,width=3.5pt]}]
        (P) -- ++(\Aer+180:15.5mm);
  \node[note, guid, anchor=south east, text width=25mm, align=right]
       at ($(P)+(\Aer+180:16.5mm)+(0,1.5)$)
       {$\dot{r}$ — $\hat{e}_r$ 축의 그림자. 지금은 \textbf{뒤쪽}이라 음수};

  \draw[thr, line width=1.2pt, -{Stealth[length=4.5pt,width=3.5pt]}]
        (P) -- ++(\Aet:22.1mm);
  \node[note, thr, anchor=north west, text width=25mm] at ($(P)+(\Aet:17mm)+(1.5,-1.5)$)
       {$r\dot{\theta}$ — $\hat{e}_\theta$ 축의 그림자};

  \node[note, anchor=north west, text width=54mm] at (-30,-14)
       {$\dot{r}$ 은 \textbf{중심까지 거리가 줄어드는 빠르기},
        $r\dot{\theta}$ 는 \textbf{중심 둘레를 도는 빠르기}다.
        두 축이 직각이므로 두 그림자를 다시 합치면 원래 화살표가 된다};

  \node[wnote, anchor=north west, text width=54mm] at (-30,-33)
       {$\dot{\theta}$ 가 아니라 $r\dot{\theta}$ 인 까닭 — $\dot{\theta}$ 는 각속도
        [rad/s] 라 속도가 아니다. 호 길이 $=$ 반경 $\times$ 각 이므로
        $r$ 을 곱해야 [m/s] 가 된다. 같은 각속도라도 바깥에 있을수록 빨리 돌아야 한다};

\end{scope}

\end{tikzpicture}
\end{document}
